\section{Related Work}
\label{sec:related}

\noindent\textbf{Speech-Driven 3D Facial Animation.}\quad
Speech-driven 3D facial animation generates temporally synchronized facial motion from audio, and existing methods are built as direct regression from acoustic features to facial geometry. Early efforts predict vertex displacements, where VOCA~\citep{cudeiro2019voca} produces speaker-specific vertices and MeshTalk~\citep{richard2021meshtalk} learns a categorical latent space covering non-verbal motion, and Transformer-based models later improved long-range temporal modeling~\citep{fan2022faceformer,fan2024unitalker}. Because the audio-to-motion mapping is one-to-many, these regressors average over plausible motions and produce over-smoothed results, which CodeTalker~\citep{xing2023codetalker} mitigates with discrete motion priors and which diffusion-based methods~\citep{park2023said,stan2023facediffuser,sun2025vividtalk} address by modeling motion stochastically. In production settings, Audio2Face-3D~\citep{chung2025audio2face} and related avatar pipelines~\citep{wei2024aniportrait,lin2025cyberhost,sun2024uniavatar,zhen2025teller} provide real-time animation, SentiAvatar~\citep{jin2026sentiavatar} plans at the sentence level, and Ex-Omni~\citep{zhang2026ex} uses an omni-modal model for semantic understanding while delegating coefficient prediction to a lightweight generator. Across this line, the linguistic structure of speech never enters the module that decides mouth shape.

\noindent\textbf{Linguistic and Articulatory Priors.}\quad
Speech articulation is tied to phonetic structure, where bilabial consonants require lip closure, open vowels require larger jaw opening, and phoneme transitions shape the temporal profile of motion~\citep{cohen1993modeling,taylor2012dynamic,edwards2016jali}. Viseme-based methods encode these relations as explicit rules or discrete inventories~\citep{zhou2018visemenet,taylor2017deep,bao2023learning}, which makes articulation interpretable at the cost of collapsing continuous coarticulation into a few categories. Later methods relax this constraint by supplying higher-level cues as auxiliary conditions, including emotion signals~\citep{peng2023emotalk,danvevcek2023emotional,gan2023efficient,jin2026sentiavatar} and linguistic features~\citep{fan2022joint,xu2024kmtalk,wu2025keyframeface}. In these designs the cue enters as a feature vector, so the model still has to learn from scratch what facial action each phonetic category names.

\noindent\textbf{Positioning of AudioMouth.}\quad
The output representation decides whether such grounding is even expressible. Vertex-based representations~\citep{cudeiro2019voca,fan2022faceformer,aloni2025express4d} preserve fine detail but carry no semantics, and parametric models such as FLAME~\citep{li2017learning} gain compactness through PCA axes that are statistical instead of articulatory, so in both cases a prediction cannot be named. ARKit blendshapes~\citep{apple_arkit} instead define named facial actions, which lets one vocabulary describe the articulatory cause and the generated control. In this work, we focus on the articulatory structure of speech, with the goal of making it an explicit part of both the input and the output of mouth-motion generation. Unlike previous approaches that treat linguistic information as auxiliary features of an acoustic regressor, AudioMouth adopts a more deliberate design by generating named ARKit controls inside a multimodal language model that is told which articulatory action each phoneme group requires, and by evaluating the decoded trajectory as a whole through a reward whose geometric term reads out how the predicted coefficients deform the mouth surface.
